\documentclass[10pt]{article}

\usepackage[utf8]{inputenc}

\usepackage[T1]{fontenc}

\usepackage{amsmath}

\usepackage{amsfonts}

\usepackage{amssymb}

\usepackage[version=4]{mhchem}

\usepackage{stmaryrd}

\usepackage{bbold}



\begin{document}

Находят приложение в задачах математич. физики В. п., не являющиеся однородными, в К-рых учитывается «возраст» частиц, а также немарковские В. П. (см. Яноши уравнение).



Лum.: [1] Ба руч а - Р и д А. Т., Элементы теории марковских процессов и их приложения, пер. с англ., М., 1969; [2] Се в астьян о в Б. А., Ветвящиеся процессы, М., 1971. Ю. П. Вирченко.



ВЕЩЕСТВЕННАЯ АЛГЕБРА ЛИ - см. ЛИ алгебра. ВЕЩЕСТВЕННАЯ АЛГЕБРАИЧЕСКАЯ КРИВАЯ - $\mathrm{cM}$. Клейнова поверхность.



ВЕЩЕСТВЕННАЯ АНАЛИТИЧЕСКАЯ ФУНКЦИЯ - см. Клейнова поверхность.



ВЕЩЕСТВЕННОЕ СУПЕРПРОСТРАНСТВО - см. Cyперпространство.



ВЕЩЕСТВЕННЫЙ АНАЛИТИЧЕСКИЙ ФУНКЦИОНАЛ - см. Аналитический функционал.



ВЗАИМОДЕЙСТВИЕ в ф и з и е - воздействие тел или частиц друг на друга, приводящее к изменению состояния их движения. В механике Ньютона взаимное действие тел друг на друга характеризуется силой. Более обшей характеристикой В. является потенциальная энергия.



Первоначально в физике утвердилось представление о том, что В. между телами может осуществляться непосредственно через пустое пространство, к-рое не принимает никакого участия в передаче В.; при этом передача В. происходит мгновенно. Так, считалось, что перемещение Земли должно сразу же приводить к изменению силы тяготения, действующей на Луну. В этом состояла так наз. концепция даль ь о де йс т в и я.



Однако эти представления были оставлены как не соответствующие действительности после открытия и исследования электромагнитного поля. Было доказано, что В. электрически заряженных тел осуществляется не мгновенно и перемещение одной заряженной частицы приводит к изменению сил, действующих на другие частицы, не в тот же момент, а спустя конечное время. В разделяющем частицы пространстве происходит нек-рый процесс, к-рый распространяется с конечной скоростью. Соответственно имеется «посредник», осуществляюший В. между заряженными частицами. Этот посредник был назван эле к тр о магн и т ны м п олем. Каждая электрически заряженная частица создает электромагнитное поле, действующее на другие частицы. Скорость распространения электромагнитного поля равна скорости света в вакууме $c \approx 3 \cdot 10^{10} \mathrm{~cm} / \mathrm{c}$. Возникла новая концепция - бл и з ко де й с т в и я, к-рая позже была распространена и на любые другие В. Согласно этой концепции, В. между телами осуществляется посредством тех или иных полей, непрерывно распределенных в пространстве. Так, всемирное тяготение осуществляется гравитационным полем.



После появления квантовой теории поля представление о В. существенно изменилось. Согласно этой теории, любое поле представляет собой совокупность частиц - квантов этого поля. Каждому полю соответствуют свои частицы. Напр., квантами электромагнитного поля являются фото$н ы$, то есть фотоны являются переносчиками этого В. Аналогично другие виды В. возникают в результате обмена между частицами квантами соответствующих полей.



Несмотря на разнообразие воздействий тел друг на друга (зависящих от В. слагающих их элементарных частиц), в природе, по современным данным, имеется лишь 4 типа фундаментальных В. Это (в порядке возрастания интенсивности В.): гравитационное взаимодействие, слабое взаимодействие (отвечающее за большинство распадов и многие превращения элементарных частиц), электромагнитное взаимодействие, сильное взаимодействие (обеспечивающее, в частности, связь частиц в атомных ядрах). Интенсивность В. определяется соответствующей константой взаимодействия, или константой связи. В частности, для электромагнитного В. константой связи является электрич. заряд. Квантовая теория электромагнитного В.- квантовая электродинамика - превосходно описывает все известные электромагнитные явления. Слабое В. осуществляется посредством промежуточных векторных бозонов. Найдена глубокая связь слабого В. с электромагнитным, что привело к их объединению в элект- рослабое взаимодействие. Основу сильного В., по современным представлениям, составляет В. между составными частями адронов - кварками. Это В., переносчиками к-рого служат глюоны, определяется особой константой взаимодействия - цветом и описывается квантовой хромодинамикой. В. адронов друг с другом представляет собой лишь остаточный эффект межкварковых сил, подобно тому как молекулярные силы - остаточный эффект кулоновского В. электронов и ядер молекул. Делаются попытки объединения слабого, электромагнитного и сильного В. (так наз. великого объединения модели), а также всех видов В., включая гравитационное (см. Супергравитация).



Г.Я. Мякишев.



ВЗАИМОДЕЙСТВИЯ КОНСТАНТА - см. ВЗаимодействие, Связи константа.



ВЗАИМОДЕЙСТВИЯ ОБЛАСТЬ - область $\Omega_{n}(t)$ в конфигурационном пространстве классической статистической системы $n$ частиц $\left(\mathbb{R}^{v n}, v=1,2,3\right)$ такая, что если координаты частиц $\left(x_{1}, \ldots, x_{n}\right) \in \Omega_{n}(t) \subset \mathbb{R}^{v n}$, то все $n$ частиц могут провзаимодействовать друг с другом за время $t$, то есть система $n$ частиц может существовать как единое целое (объем В. о. зависит от начальных импульсов частиц $p_{1}, \ldots$, $\left.p_{n}\right)$. Напротив, если $\left(x_{1}, \ldots, x_{n}\right) \subset \mathbb{R}^{v n} \backslash \Omega_{n}(t)$, то система распадается как минимум на 2 группы частиц, невзаимодействуюших друг с другом на отрезке времени $[0, t]$. Понятие В. о. введено Д. Я. Петриной [1].



Для одномерных $(v=1)$ систем частиц с парным финитным потенциалом взаимодействия $\Phi(|x|) \in C^{2}(2 a, \infty)$ с твердой сердцевиной (радиуса $a>0$ ) средний по максвелловскому распределению объем В.о. можно оценить сверху так $(A, B, C, D<\infty-$ постоянные, зависящие от потенциала взаимодействия):



\[

V_{\mathrm{cp}}\left(\Omega_{n}(t)\right) \leqslant A^{n}\left(B t^{2}+C t+D\right)^{\mathrm{vn}} \cdot n^{n} .

\]



В случае $v>1$ оценка, аналогичная $\left({ }^{*}\right)$, получена с использованием дополнительных ограничений технич. характера на потенциал и допустимые начальные данные (см. [3]).



Использование оценки (*) позволяет доказать существование термодинамич. предела для последовательности функций распределения, удовлетворяющих Боголюбова уравнениям.



Лит.: [1] Пе т р и н а Д. Я., «Теоретич. и математич. физика», 1979, т. 38, с. 230; [2] ПетринаД. Я., ГерасименкоВ.И., Малыш е в П. В., Математические основы классической статистической механики, Киев, 1985; [3] Petrina D.Ja., Ge rasimenko V.I., Malyshev P. V., «Sov. Sci. Rev., C, Math. Phys. Rev.», 1988, v. 7, p. 281 .



П. В. Малышев.



ВЗАИМОДЕЙСТВИЯ ПОТЕНЦИАЛ - скалЯрная энергетическая характеристика воздействия силового поля на внесенную в него частицу (тело). Для частицы в электростатич. (гравитационном) поле В. п. равен отношению потенциальной энергии взаимодействия частицы с полем к величине ее заряда (массы). Иными словами, В. п. с полем для частицы с единичным зарядом (массой) равен работе, к-рую необходимо затратить, чтобы переместить частицу из точки, потенциал к-рой условно принимается равным нулю, в данную точку поля.



Потенциальная энергия системы $\boldsymbol{n}$ частиц в общем случае определяется так:



\[

U\left(x_{1}, \ldots, x_{n}\right)=\sum_{k=1}^{n} \sum_{1 \leqslant i_{1}<\ldots<i_{k} \leqslant n} \Phi^{(k)}\left(x_{i_{1}}, \ldots, x_{i_{k}}\right),

\]



где $x_{i} \in \mathbb{R}^{v}, v=1,2,3, i \in\{1, \ldots, n\}-$ координаты частиц, а $\Phi^{k}-k$-частичные В. п. В большинстве случаев предполагается, что все В. П., кроме парных $(k=2)$, равны нулю, но в нек-рых задачах используются и многочастичные потенциалы $(k>2), \Phi^{(1)}$ обычно имеет смысл В. П. с внешним полем. Потенциальная энергия системы частиц с парным В. п. имеет вид:



\[

U\left(x_{1}, \ldots, x_{n}\right)=\sum_{1 \leqslant i<j \leqslant n} \Phi\left(x_{i}, x_{j}\right)

\]





\end{document}